\documentclass[10pt,twocolumn,twoside,a4paper]{report} % article

\usepackage[slovak,english]{babel}
\usepackage[IL2]{fontenc}
\usepackage[utf8]{inputenc}
\usepackage{graphicx}
\usepackage{url}
\usepackage{hyperref}

\usepackage{cite}

\graphicspath{ {./images/} }

\pagestyle{headings}

\title{Personalized Recommendation Systems in Social Media\thanks{Semester project for the course Engineering Methods, academic year 2024/25, supervised by Ing. Richard Marko, PhD.}} % name and title of the course supervisor

\author{Yevhenii Deshchenko\\[2pt]
	{\small Slovak Technical University in Bratislava}\\
	{\small Faculty of Informatics and Information Technology}\\
	{\small \texttt{xdeshchenko@stuba.sk}}
	}

\date{\small 10 October 2024}

\begin{document}

\maketitle

\begin{abstract}
With the growing popularity of social media, recommendation systems have become increasingly complex and sophisticated. 
These systems rely on a variety of algorithms, including those based on graph theory, statistics, machine and deep learning. 
Content-based recommendation systems vary significantly across many different social media platforms, as they depend on the 
type of primary content users share and the options provided for user interaction with the content. 
Most approaches use as much data as available, leading to massive data collection that raises multiple concerns, such as privacy issues.


As social media platforms continue to grow in popularity, the inner workings of these algorithms are often kept secret,
as they have become valuable tools for third parties seeking to exploit them for visibility and influence.
This exploitation frequently aims to push content with political or other agendas, in an effort to gain popularity on these platforms.
However, there are exceptions, such as Twitter, which made its recommendation algorithm public after a data breach.
This decision led to significant criticism from its user base. The secrecy and potential for manipulation continue to
spark debates about the role and transparency of recommendation systems in shaping user experiences and content visibility on social media.
\end{abstract}

% \includegraphics[width=\textwidth]{fiit.eps}

\section{Introduction}

In a modern era where social media has become an integral part of daily life, personalized recommendation systems play a
crucial role in enhancing marketing, entertainment, and the overall user experience. By predicting and delivering content
that resonates with users, these systems contribute significantly to social media growth, fostering higher user engagement
and supporting an interconnected ecosystem of content creators, businesses, and users alike.

While these systems are built upon statistical and mathematical foundations, practical implementations vary widely.
They often combine several approaches, including complex algorithms, databases, big data analytics, machine learning,
and deep learning techniques, to optimize predictions and deliver tailored content. 

However, the reliance on extensive user data raises privacy concerns, as large amounts of personal information must
be stored and processed. Additionally, biases within these algorithms can lead to unintended consequences, potentially
skewing recommendations or enabling exploitation. This paper will explore the methodologies behind recommendation systems,
their societal impacts, and the balance between personalization and ethical considerations.

\section{Recommendation systems overview} \label{rec_sys_overview}
The primary goal of recommendation systems (RS) is to suggest content, connections, posts, or advertisements that may interest the end user.
Social media, as a technology, has evolved over time and is now embedded into blogs, forums, and other digital forms,
facilitating community building, content sharing, and information exchange \cite{anandhan2018social}.
These systems have become integral in helping users navigate vast amounts of content efficiently, allowing platforms to deliver more personalized experiences.

\subsection{Social Media overview}
Social media platforms can be classified into few categories \cite{anandhan2018social} based on their primary function and user interaction patterns:
\begin{itemize}
	\item \textbf{Digital Libraries:} Platforms like Pinterest and Flickr act as digital libraries where users can share, categorize, and access large collections of images, documents, or other resources. Recommendation systems on these platforms focus on content discovery, suggesting relevant media or documents based on specific determined user preferences.
    \item \textbf{Forums:} Platforms such as Reddit and Stack Overflow allow users to engage in topic-focused discussions and information sharing. Here, recommendation systems suggest relevant threads, discussions, or posts based on user interactions and interest tags.
    \item \textbf{Blogs:} Platforms like Medium and Tumblr are designed for content creation, where users can write articles, stories, or posts. Recommendation systems in blogs aim to match readers with articles that align with their interests, reading history, and trending topics.
    \item \textbf{Social Networks:} Platforms such as Facebook, Instagram, and Twitter are built around user-to-user connections and content sharing. These platforms heavily rely on recommendation systems to suggest new connections, posts, and advertisements, fostering interaction between content creators, advertisers, and consumers.
    \item \textbf{Other:} Platforms, like YouTube, which blend in themselves several of the above categories.
\end{itemize}

One of the most widespread applications of recommendation systems is within social networks, such as Facebook, Instagram, or Twitter.
These systems play a pivotal role in connecting content creators and advertisers with consumers by suggesting relevant posts, trending topics,
or advertisements tailored to user preferences.
As users engage with these platforms, recommendation systems analyze behavior, such as likes, shares, comments, and even search queries,
to curate personalized content. This personalization aims to enhance user engagement, increase time spent on the currentplatform,
as well as improve overall user satisfaction.

\subsection{Types of recommendation systems}
Algorithms behind personal recommendation systems are typically classified into either content-based filtering or collaborative filtering. \cite{raghavendra2018personalized}

In content-based filtering, algorithm must define or in one way or another construct a certain profile of the target user. Later, based on that
built profile the algorithm decides what items, whether it would be posts or advertisements, would be approppreate to recommend.
The way it works, is that it firstly considers item characteristics defined and assignes them keywords and tags. Using machine learning algorithms,
decision trees and clustering it recommends products with similar features to those previously liked or positively reviewed. \cite{raghavendra2018personalized}\cite{contentvscollaborative}

In collaborative filtering, recommendation system acts based on user's behavioural analysis. 

\section{Data collection and privacy issues} \label{data_issues}
\section{Content manipulation and transparency challanges} \label{content_manipulation}
\section{Future challanges in recommendation systems} \label{content_manipulation}

\section{Conclusion} \label{conclusion}

%\acknowledgement{If you want to thank someone\ldots}

% \nocite{*}
\bibliography{literatura}
\bibliographystyle{plain} % or alpha, abbrv, or any other style
\end{document}
